% !TeX root = ../cosmology_summary.tex

%% --------------------------------------------------------
\chapter{Розігрів і народження матерії}
\label{sec:reheating}
%% --------------------------------------------------------

Коли інфлатон досягає мінімуму потенціалу $V(\phi)$, умова
інфляції порушується ($\varepsilon \to 1$) і прискорене
розширення припиняється. Але Всесвіт у цей момент
перебуває у парадоксальному стані: він розширився в
$e^{50}$--$e^{60}$ разів, охолов майже до абсолютного нуля
і є практично порожнім~--- вся енергія зосереджена у
вигляді однорідних класичних коливань поля $\phi$ навколо
мінімуму. Жодної плазми, жодних часток Стандартної моделі.

Перехід від цього «холодного» стану до гарячого
ультрарелятивістського Всесвіту, з якого починається
стандартна космологія Великого вибуху, описується теорією
\emph{розігріву} \cite{linde_1982, albrecht_1982}. Динаміку
розігріву поділяють на дві послідовні стадії: спочатку
нелінійний вибуховий \textbf{пре-розігрів}, де параметричний
резонанс за лічені осциляції перекачує енергію з поля у
частинки, а потім повільний \textbf{пертурбативний розпад},
де залишки інфлатона розпадаються і система термалізується.

%%% --------------------------------------------------------
%\section{Прерозігрів та параметричний резонанс}
%%% --------------------------------------------------------
%
%Поблизу мінімуму параболічного потенціалу
%$V(\phi) = \frac{1}{2}m^2\phi^2$, рівняння \eqref{eq:klein_gordon} має вигляд:
%\begin{equation*}
%     \ddot{\phi} + 3H\dot{\phi} + m^2\phi = 0.
%\end{equation*}
%Оскільки частота масових осциляцій значно перевищує швидкість розширення Всесвіту ($m_{\phi} \gg H$), у WKB-наближенні повільно мінливої амплітуди розв'язок шукається у вигляді
%\begin{equation*}
%    \phi(t) \approx \Phi_0(t)\sin(mt),
%\end{equation*}
%що після підстановки та нехтування вищими похідними ($\dot{\Phi}_0 \ll m\Phi_0$) зводить систему до диференціального рівняння першого порядку $\dot{\Phi}_0 + \frac{3}{2}H\Phi_0 = 0$. Інтегрування цього виразу дає степеневий закон загасання амплітуди $\Phi_0(t) \propto a(t)^{-3/2}$, внаслідок чого усереднена за період коливань густина енергії поля падає як $\rho_\phi \approx \frac{1}{2}m^2\Phi_0^2 \propto a^{-3}$, змушуючи квантовий конденсат імітувати термодинамічну поведінку звичайної пилоподібної матерії.
%
%%Поблизу мінімуму параболічного потенціалу
%%$V(\phi) = \frac{1}{2}m^2\phi^2$ інфлатон осцилює як
%%$\phi(t) \approx \Phi_0(t)\sin(mt)$, де амплітуда
%%$\Phi_0(t)$ повільно затухає через хабблівське тертя.
%Якщо інфлатон пов'язаний з іншим скалярним полем $\chi$
%через член взаємодії $g^2\phi^2\chi^2$, то ефективна маса
%поля $\chi$ осцилює в часі: $m_\chi^2(t) \supset
%g^2\Phi_0^2\sin^2(mt)$. Рівняння руху для $k$-ї моди
%$\chi_k$ набуває форми рівняння Матьє
%\cite{kofman_1994, kofman_1997}:
%\begin{equation}\label{eq:mathieu}
%  \ddot{\chi}_k + 3H\dot{\chi}_k
%  + \bigl[k_\mathrm{phys}^2 + m_\chi^2
%         + g^2\Phi_0^2\sin^2(mt)\bigr]\chi_k = 0.
%\end{equation}
%
%Рівняння з періодично змінними коефіцієнтами має зони
%нестабільності, де амплітуда розв'язку зростає
%експоненціально~--- \emph{параметричний резонанс}.
%Народження часток $\chi$ відбувається не поштучно, а
%вибухово: бозонні поля мають стимульоване народження у
%вже заповнені стани (аналог ефекту Бозе-Айнштайна), тому
%заселеність кожної резонансної моди зростає за лічені
%осциляції інфлатона. Цей процес~--- \textbf{пре-розігрів}
%(preheating)~--- за короткий час перекачує значну частку
%енергії поля у частинки $\chi$.
%% --------------------------------------------------------
\section{Прерозігрів та параметричний резонанс}
%% --------------------------------------------------------

Поблизу мінімуму параболічного потенціалу $V(\phi) = \frac{1}{2}m^2\phi^2$, рівняння \eqref{eq:klein_gordon} має вигляд:
\begin{equation*}
     \ddot{\phi} + 3H\dot{\phi} + m^2\phi = 0.
\end{equation*}
Оскільки частота масових осциляцій значно перевищує швидкість розширення Всесвіту ($m \gg H$), у WKB-наближенні повільно мінливої амплітуди розв'язок шукається у вигляді
\begin{equation*}
    \phi(t) \approx \Phi_0(t)\sin(mt),
\end{equation*}
що після підстановки та нехтування вищими похідними ($\dot{\Phi}_0 \ll m\Phi_0$) зводить систему до диференціального рівняння першого порядку $\dot{\Phi}_0 + \frac{3}{2}H\Phi_0 = 0$. Інтегрування цього виразу дає степеневий закон загасання амплітуди $\Phi_0(t) \propto a(t)^{-3/2}$, внаслідок чего усереднена за період коливань густина енергії поля падає як $\rho_\phi \approx \frac{1}{2}m^2\Phi_0^2 \propto a^{-3}$, змушуючи квантовий конденсат імітувати термодинамічну поведінку звичайної пилоподібної матерії.

Якщо інфлатон пов'язаний з іншим квантованим скалярним полем $\chi$ через лагранжіан взаємодії $\mathcal{L}_{\mathrm{int}} = -\frac{1}{2}g^2\phi^2\chi^2$, то коливання інфлатона створюють періодичний зовнішній фон. Ефективна маса поля $\chi$ починає осцилювати в часі: $m_{\mathrm{eff}}^2(t) = m_\chi^2 + g^2\Phi_0^2\sin^2(mt)$. Рівняння руху для $k$-ї Фур'є-моди $\chi_k$ набуває форми рівняння Матьє у просторі, що розширюється \cite{kofman_1994, kofman_1997}:
\begin{equation}\label{eq:mathieu}
  \ddot{\chi}_k + 3H\dot{\chi}_k + \bigl[k_\mathrm{phys}^2 + m_\chi^2 + g^2\Phi_0^2\sin^2(mt)\bigr]\chi_k = 0.
\end{equation}

Щоб зрозуміти, чому саме рівняння Матьє керує народженням часток, корисно провести аналогію з механікою: рівняння~\eqref{eq:mathieu} є точним аналогом рівняння маятника зі змінною довжиною підвісу. Якщо ритмічно смикати нитку маятника з подвоєною власною частотою, розгойдування наростатиме експоненціально~--- навіть від нескінченно малого початкового відхилення. У нашому випадку роль «підвісу» відіграє ефективна маса поля $\chi$: коливання інфлатона periodично змінюють $m_{\rm eff}^2(t)$, а отже й «жорсткість пружини» в рівнянні руху для кожної моди $\chi_k$. Принципова різниця від механічного маятника~--- квантова: навіть якщо поле $\chi$ на початку перебуває у вакуумному стані ($N_k = 0$), нульові коливання забезпечують ненульову «початкову амплітуду», від якої резонанс починає розгойдуватися. Саме тому параметричний резонанс є неминучим за наявності будь-якої зв'язки $g \neq 0$~--- він не потребує «посівного» теплового фону.

Для аналізу нестабільності зручно зробити заміну поля $\chi_k = X_k \cdot a^{-3/2}$ та перейти до безрозмірного часу $z = mt$, що усуває хабблівське тертя і зводить вираз до канонічного вигляду $\frac{d^2 X_k}{dz^2} + [A_k - 2q\cos(2z)]X_k = 0$, де параметр резонансу дорівнює $q = \frac{g^2\Phi_0^2}{4m^2}$. Це диференціальне рівняння з періодичними коефіцієнтами має чітко визначені зони нестабільності на площині параметрів $(A_k, q)$. Коли частота фонових коливань інфлатона потрапляє в резонансну смугу, амплітуда моди $X_k$ починає зростати експоненціально: $X_k \propto e^{\mu_k mt}$, де $\mu_k$ --- показник нестабільності.

На мові квантової теорії поля цей процес описується як \emph{параметричний резонанс}, що призводить до неадіабатичної еволюції вакуумного стану. Народження часток $\chi$ визначається через коефіцієнти Боголюбова $\beta_k$, а чисельність заповнення квантових станів росте як $N_k(t) = |\beta_k|^2 \propto e^{2\mu_k mt}$. Таке народження відбувається не поштучно, а вибухово: бозонні поля мають стимульоване народження у вже заповнені стани (що є прямим наслідком квантової статистики Бозе---Айнштайна), тому заселеність кожної резонансної моди подвоюється за лічені осциляції інфлатона. Цей початковий етап --- \textbf{пре-розігрів} (preheating) --- є суттєво нелінійним і за надкороткий час перекачує значну частку кінетичної енергії однорідного конденсату інфлатона у високоенергетичні кванти поля $\chi$, закладаючи фундамент для подальшої повної термалізації Всесвіту.

%%% --------------------------------------------------------
%\section{Пертурбативний розпад та термалізація}
%%% --------------------------------------------------------
%
%Коли амплітуда $\Phi_0$ падає нижче порогу резонансу,
%параметричний резонанс вимикається. Подальший розпад
%залишків інфлатона відбувається повільно і описується
%пертурбативною теорією: у рівняння збереження енергії
%вводиться феноменологічна ширина розпаду $\Gamma_\phi$:
%\begin{align}
%  \dot{\rho}_\phi + 3H\rho_\phi
%    &= -\Gamma_\phi\,\rho_\phi,
%    \label{eq:reheat_rho_phi} \\
%  \dot{\rho}_r + 4H\rho_r
%    &= +\Gamma_\phi\,\rho_\phi.
%    \label{eq:reheat_rho_r}
%\end{align}
%Перше рівняння~--- розпад інфлатона; друге~--- накопичення
%радіації. Коефіцієнти $3H$ і $4H$ відображають розбавлення
%густин розширенням: $\rho_\phi \propto a^{-3}$ (матерія),
%$\rho_r \propto a^{-4}$ (радіація).
%
%Поки $H \gg \Gamma_\phi$, розширення домінує і розпад
%відбувається повільно. Процес завершується при
%$H \sim \Gamma_\phi$: залишки інфлатона розпадаються майже
%миттєво, а релятивістські продукти термалізуються через
%пружне і непружне розсіяння, встановлюючи єдину температуру.
%
%\begin{keybox}{Температура розігріву $T_\mathrm{rh}$}
%визначається з умови $H^2 = \frac{8\pi G}{3}\rho_r$ при
%$H \approx \Gamma_\phi$ і рівноважного розподілу радіації
%$\rho_r = \frac{\pi^2}{30}g_* T^4$:
%\begin{equation}\label{eq:t_reheating}
%  T_\mathrm{rh}
%    = \left(\frac{90}{8\pi^3 g_*}\right)^{\!1/4}
%      \sqrt{M_\mathrm{Pl}\,\Gamma_\phi},
%\end{equation}
%де $g_* \sim 100$~--- ефективна кількість
%релятивістських ступенів вільності Стандартної моделі.
%\end{keybox}
%
%Температура $T_\mathrm{rh}$ задає початкові умови для
%всієї подальшої термічної історії: баріогенезису, фазових
%переходів і утворення темної матерії. Для нашої задачі
%важливо інше: саме в момент термалізації гаряча
%фотон-баріонна плазма «успадковує» просторові
%неоднорідності кривини $\mathcal{R}_k$, заморожені під час
%інфляції. З цього моменту вони більше не є статичними~---
%тиск фотонного газу змушує плазму коливатися, і саме ці
%акустичні коливання ми бачимо сьогодні у вигляді
%температурних плям на карті CMB (рис.~\ref{fig:planck_map})
%і піків на її кутовому спектрі (рис.~\ref{fig:cmb_tt}).

%% --------------------------------------------------------
\section{Пертурбативний розпад та термалізація}
%% --------------------------------------------------------

Коли амплітуда коливань інфлатона падає нижче критичного порогу $\Phi_0 \ll \Phi_{\mathrm{crit}}$, нелінійні ефекти неадіабатичного пре-розігріву затухають, і параметричний резонанс повністю вимикається. Подальша дисипація залишків енергії конденсату однорідного поля інфлатона відбувається повільно на квантовому рівні через поштучний розпад на легкі ступені вільності Стандартної моделі. Ця стадія описується пертурбативною теорією розпадів, де мікрофізичні процеси враховуються через феноменологічну ширину розпаду $\Gamma_\phi$, обчислену за золотим правилом Фермі для відповідних діаграм Фейнмана.

З погляду класичної макроскопічної фізики, цей процес є прямим аналогом дисипації механічної енергії зануреного у рідину маятника за рахунок в'язкого тертя Ньютона, де робота сил опору безповоротно переходить у тепло. Еволюція густин енергії системи задається системою зв'язаних диференціальних рівнянь балансу, які узагальнюють класичний закон радіоактивного розпаду Резерфорда---Содді для двох послідовних резервуарів:
\begin{align}
  \dot{\rho}_\phi + 3H\rho_\phi
    &= -\Gamma_\phi\,\rho_\phi,
    \label{eq:reheat_rho_phi} \\
  \dot{\rho}_r + 4H\rho_r
    &= +\Gamma_\phi\,\rho_\phi.
    \label{eq:reheat_rho_r}
\end{align}
Рівняння \eqref{eq:reheat_rho_phi} описує експоненціальне квантове зменшення енергії інфлатона через канал розпаду, тоді як рівняння \eqref{eq:reheat_rho_r} задає динаміку джерела для утворення вторинної ультрарелятивістської радіації (газу частинок, чия кінетична енергія значно перевищує їхню масу спокою, $E \approx pc \gg m_0 c^2$).

Доданки $3H\rho_\phi$ та $4H\rho_r$ у лівих частинах є чисто класичними гідродинамічними наслідками~\eqref{eq:conservation_law}.

Математична розбіжність у коефіцієнтах строго визначається відповідними рівняннями стану ($p = w\rho$):
\begin{itemize}
    \item Для нерелятивістського конденсату інфлатона тиск $p_\phi \approx 0$ (рівняння стану пилу $w=0$), що після підстановки $\dot{\rho}_\phi + 3H(\rho_\phi + 0) = 0$ дає класичне збереження маси у супутньому об'ємі: $\rho_\phi \propto a^{-3}$. Енергія розбавляється суто через геометричне збільшення об'єму простору $V \propto a^3$.
    \item Для ультрарелятивістського газу продуктів розпаду тиск становить $p_r = \frac{1}{3}\rho_r$ (рівняння стану радіації $w=1/3$). Це приводить до виразу $\dot{\rho}_r + 3H\left(\rho_r + \frac{1}{3}\rho_r\right) = \dot{\rho}_r + 4H\rho_r = 0$, звідки випливає закон розбавлення $\rho_r \propto a^{-4}$. Фізично це означає, що окрім тривимірного геометричного розрідження ($a^{-3}$), релятивістські частинки додатково втрачають енергію ($a^{-1}$) через доплерівське червоне зміщення внаслідок розширення самої тканини простору.
\end{itemize}

Поки темп розширення значно перевищує швидкість мікрофізичних процесів ($H \gg \Gamma_\phi$), хабблівське розрідження домінує, і густина енергії випромінювання перебуває у динамічній рівновазі між припливом енергії від розпаду та її охолодженням розширенням, поводячись як $\rho_r \propto a^{-3/2}$. Пертурбативний розпад завершується, коли параметри зрівнюються: $H \approx \Gamma_\phi$. У цей момент залишки конденсату інфлатона розпадаються майже миттєво. Первинні високоенергетичні продукти розпаду, що перебували у сильно нерівноважному стані, починають інтенсивно взаємодіяти.

Під \textbf{термалізацією} в космології розуміють макроскопічний фазовий перехід системи від сильно нерівноважного кінетичного стану до стану локальної термодинамічної рівноваги, що супроводжується максимізацією ентропії. Цей процес відбувається через каскад пружних і непружних (із незбереженням кількості частинок) процесів розсіяння. Математично цей етап описується класичним кінетичним рівнянням Больцмана, де інтеграл зіткнень швидко перерозподіляє енергію по ступенях вільності. Це призводить до затирання будь-якої пам'яті про початковий спектр продуктів розпаду та встановлює розподіл Планка з єдиною температурою для всіх релятивістських компонентів плазми.

\begin{keybox}{Температура розігріву $T_\mathrm{rh}$}
Температура розігріву визначається з умови рівності параметра Габбла і ширини розпаду $H \approx \Gamma_\phi$. Підставляючи класичний закон Стефана---Больцмана для густини енергії чорного тіла $\rho_r = \frac{\pi^2}{30}g_* T^4$ у перше рівняння Фрідмана $H^2 = \frac{8\pi G}{3}\rho_r = \frac{8\pi}{3 M_{\mathrm{Pl}}^2}\rho_r$, отримуємо фінальне значення температури завершення розігріву:
\begin{equation}\label{eq:t_reheating}
  T_\mathrm{rh}
    = \left(\frac{90}{8\pi^3 g_*}\right)^{\!1/4}
      \sqrt{M_\mathrm{Pl}\,\Gamma_\phi},
\end{equation}
де $M_{\mathrm{Pl}} = G^{-1/2} \approx 1.22 \times 10^{19}$~GeV --- планківська маса, а $g_* \sim 100$ --- ефективна кількість релятивістських ступенів вільності Стандартної моделі при високих енергіях.
\end{keybox}

Величина $T_\mathrm{rh}$ є фундаментальним космологічним параметром, що задає початкові термодинамічні умови для всієї подальшої гарячої історії Всесвіту: від баріогенезису та електрослабкого фазового переходу до загартування первинного нуклеосинтезу (BBN) і термодинамічного відщеплення темної матерії.

Для нашої задачі критично важливим є наступний перехід: саме в момент остаточної термалізації утворена гаряча фотон-баріонна плазма дзеркально \enquote{успадковує} просторовий рельєф первинних неоднорідностей кривини $\mathcal{R}_k$, які були згенеровані квантовими флуктуаціями та заморожені поза космологічним горизонтом під час інфляції. Щойно радіація стає панівним компонентом Всесвіту, ці супергоризонтні збурення починають входити під горизонт. З цього моменту вони втрачають свій статичний характер --- колосальний градієнт тиску ультрарелятивістського фотонного газу виступає як повертаюча сила, змушуючи речовину коливатися.

Цей процес повністю є еквівалентним класичній задачі гідродинаміки про поширення звуку в суцільному середовищі: градієнт тиску грає роль пружної сили, а баріони забезпечують інерційну масу. Саме ці гідродинамічні акустичні коливання, зафіксовані в епоху рекомбінації, ми спостерігаємо сьогодні у вигляді анізотропії температури на карті реліктового випромінювання (рис.~\ref{fig:planck_map}) та у формі строгої послідовності акустичних піків на її кутовому спектрі потужності (рис.~\ref{fig:cmb_tt}).


% ============================================================
% Задачі до розділу 5: Розігрів і народження матерії
\section*{Задачі}
\addcontentsline{toc}{section}{Задачі}
% ============================================================

\begin{problem}{Загасання амплітуди інфлатона під час осциляцій}
Поблизу мінімуму параболічного потенціалу $V(\phi) = \frac{1}{2}m^2\phi^2$
інфлатон осцилює за законом $\phi(t) \approx \Phi_0(t)\sin(mt)$.

\begin{enumerate}
    \item Підставте цей розв'язок у рівняння Кляйна--Ґордона
          $\ddot{\phi} + 3H\dot{\phi} + m^2\phi = 0$ і, нехтуючи $\ddot{\Phi}_0$ та $\dot{\Phi}_0^2$,
          покажіть, що $\dot{\Phi}_0 + \frac{3}{2}H\Phi_0 = 0$.
    \item Розв'яжіть це рівняння та знайдіть $\Phi_0(t)$ як функцію масштабного фактора $a(t)$.
    \item Обчисліть усереднену за період густину енергії $\rho_\phi = \frac{1}{2}m^2\Phi_0^2$
          та покажіть, що $\rho_\phi \propto a^{-3}$. Якій речовині відповідає така залежність?
\end{enumerate}
\end{problem}

% -------------------------------------------------------------------

\begin{problem}{Рівняння Матьє та параметричний резонанс}
Рівняння для моди $\chi_k$ під час пре-розігріву має вигляд
рівняння Матьє з затуханням.

\begin{enumerate}
    \item Виконайте заміну $\chi_k = X_k \cdot a^{-3/2}$ та перейдіть до безрозмірного
          часу $z = mt$. Покажіть, що рівняння \eqref{eq:mathieu} зводиться до
          \[
          \frac{d^2 X_k}{dz^2} + \bigl[A_k - 2q\cos(2z)\bigr] X_k = 0,
          \]
          де $q = \dfrac{g^2\Phi_0^2}{4m^2}$. Чому хабблівське тертя зникає після цієї заміни?
    \item Поясніть фізичний зміст параметра $q$. У якому випадку резонанс є найефективнішим?
    \item Чому народження частинок $\chi$ відбувається навіть із початково вакуумного стану
          ($N_k = 0$)? Яку роль тут відіграють нульові коливання?
\end{enumerate}
\end{problem}

% -------------------------------------------------------------------

\begin{problem}{Пертурбативний розпад та температура розігріву}
Після закінчення пре-розігріву залишки енергії інфлатона розпадаються
пертурбативно з шириною $\Gamma_\phi$.

\begin{enumerate}
    \item Використовуючи систему рівнянь \eqref{eq:reheat_rho_phi}--\eqref{eq:reheat_rho_r},
          покажіть, що в режимі $H \gg \Gamma_\phi$ густина радіації
          $\rho_r \propto a^{-3/2}$. (Підказка: припустіть, що $\rho_\phi$ домінує
          і розпадається повільно.)
    \item Умова завершення розігріву: $H \approx \Gamma_\phi$. Використовуючи
          $H^2 = \dfrac{8\pi G}{3}\rho_r$ та $\rho_r = \dfrac{\pi^2}{30}g_* T^4$,
          виведіть формулу для температури розігріву $T_{\mathrm{rh}} =
          \left(\dfrac{90}{8\pi^3 g_*}\right)^{1/4} \sqrt{M_{\mathrm{Pl}}\Gamma_\phi}$.
    \item Для типової ширини розпаду $\Gamma_\phi \sim 10^5\ \text{GeV}$ оцініть $T_{\mathrm{rh}}$.
          Чи достатньо ця температура для електрослабкого фазового переходу ($T \sim 100\ \text{GeV}$)?
\end{enumerate}
\end{problem}

% -------------------------------------------------------------------

\begin{problem}{Чи може пре-розігрів створити надлишок темної матерії?}

Параметричний резонанс може народжувати частинки $\chi$ дуже ефективно,
навіть якщо $\chi$ не має пертурбативних каналів розпаду. Припустимо,
що $\chi$ є стабільним або довгоживучим кандидатом на темну матерію.

\begin{enumerate}
    \item Оцініть максимальну густину числа $n_\chi$, яку можна отримати
          за час пре-розігріву, якщо показник нестабільності $\mu_k \sim 0.1$.
    \item Сучасне спостережне значення густини темної матерії $\Omega_{\mathrm{DM}} \approx 0.27$.
          Яку масу $m_\chi$ повинна мати частинка, щоб повністю складати темну матерію,
          народжену під час пре-розігріву?
    \item Чому такий механізм називають \emph{непертурбативним} народженням темної матерії?
          Які переваги він має перед термальним сценарієм?
\end{enumerate}
\end{problem}